\section{Purpose: distinguish magnetics, coordinates, and reconstruction}
A measured direct-axis flux curve that tilts between positive and negative
quadrature current invites several explanations: saturation, an incorrect
resistance, or a misaligned rotor frame. A curve alone cannot decide between
them. The useful question is which features an admissible magnetic map can
have, and which measurement operations introduce features outside that class.
Steady-state rotor-oriented voltage conventions follow standard machine
modeling \cite{binder2017,schroeder2021}. The companion co-energy reconstruction
framework establishes which components belong to symmetric quasi-static
machine magnetics \cite{companion-coenergy}. The present work studies the
complementary question: what information is contained in components that
violate this structure?

The central question is: given steady-state voltage, current, electrical speed,
and reconstructed flux maps, which reflection symmetries are required by a
magnetically symmetric machine, how do errors change them, and when can
resistance and angle errors be separately estimated? We keep four layers
distinct: the magnetic constitutive map, its coordinate representation,
voltage-based reconstruction, and imperfect measurements or used parameters.
Their separation matters because a coordinate error moves the argument of a
map as well as its output, whereas a resistance error changes reconstruction.

The paper proceeds from reconstruction and a compact structural foundation
to error propagation, exact frame rotation, parity residuals, and
identifiability. A normalized experiment then builds intuition from an
isotropic PM machine to a saturated cross-coupled one. It uses synthetic data,
not a claimed motor calibration or a comparison of interpolation methods.
The final sections explain correction proposals for future MeasEval analysis
and interpretation of residuals to the admissible magnetic model. Full
co-energy fitting theory and its noise/sampling experiments belong to the
companion; this paper tests error injection and parameter recovery.

\paragraph{Reading conventions.}
Unadorned quantities denote truth. Hats denote measured coordinates, measured
signals, or a used parameter, identified explicitly at each introduction.
$\flux(\current)$ is the true constitutive map; $\flux^{\rm rec}$ is a
voltage-reconstructed value. The measured map $\hat\flux(\hat\current)$
contains both frame and reconstruction effects. Differences between maps are
always compared at the same \emph{numerical current argument}. Angles are
electrical radians; speed is signed electrical angular speed. Temperature,
magnet state, and averaging convention are held fixed unless stated otherwise.
